\documentclass[11pt]{article}
\newcommand{\SheetCode}{Theory 08}
\usepackage[margin=25mm]{geometry}
\usepackage[T1]{fontenc}
\usepackage{lmodern}
\DeclareUnicodeCharacter{2605}{$\star$}
\usepackage{microtype}
\usepackage{amsmath,amssymb,mathtools,bm}
\usepackage{enumitem}
\usepackage{xcolor}
\usepackage{fancyhdr}
\usepackage{hyperref}
\usepackage[most]{tcolorbox}
\definecolor{agiBlue}{HTML}{173B57}
\definecolor{agiGold}{HTML}{B7791F}
\definecolor{agiLight}{HTML}{EFF6FA}
\definecolor{agiGreen}{HTML}{184D3B}
\hypersetup{colorlinks=true,linkcolor=agiBlue,urlcolor=agiBlue}
\setlength{\parindent}{0pt}
\setlength{\parskip}{0.55em}
\setlength{\headheight}{14pt}
\setlist{nosep,leftmargin=*}
\pagestyle{fancy}
\fancyhf{}
\fancyhead[L]{\small Part of the \href{https://github.com/mmikhasenko/GIModel.jl}{Agentic Gottried--Isgur} project}
\fancyhead[R]{\SheetCode}
\fancyfoot[C]{\thepage}
\newcommand{\dd}{\,\mathrm d}
\newcommand{\ket}[1]{\lvert #1\rangle}
\newcommand{\bra}[1]{\langle #1\rvert}
\newcommand{\braket}[2]{\langle #1\mid #2\rangle}
\newcommand{\expect}[1]{\left\langle #1\right\rangle}
\newcommand{\op}[1]{\widehat{#1}}
\newcommand{\GeV}{\,\mathrm{GeV}}
\newcommand{\R}{\mathbb{R}}
\newtcolorbox{goals}{colback=agiLight,colframe=agiBlue,title=Learning outcomes}
\newtcolorbox{reading}{colback=white,colframe=agiGold,title=Book reference,
  after skip=5mm}
\newtcolorbox{paperbridge}{colback=agiGreen!7,colframe=agiGreen,
  colbacktitle=agiGreen,coltitle=white,title=Connection to the GI paper}
\newtcolorbox{problem}[1]{colback=white,colframe=agiBlue,title={#1},
  before skip=3.5mm,after skip=1mm}
\newtcolorbox{solution}{colback=black!2,colframe=black!45,title=Solution}
\newtcolorbox{quiz}{colback=white,colframe=agiGold,title=Post-solution concept check}
\newtcolorbox{checkpoint}{colback=white,colframe=agiGold,title=Interpretation checkpoint}
\newcommand{\sheettitle}[3]{%
  \begin{center}
  {\Large\bfseries #1}\par
  \vspace{0.2em}{\color{agiBlue}#2}\par
  \vspace{0.4em}\small #3
  \end{center}\vspace{0.5em}}
\newcommand{\difficulty}[1]{\textbf{Difficulty:} #1}
\newcommand{\timebox}[1]{\textbf{Suggested time:} #1}
\newcommand{\answer}{\par\smallskip\color{black!50}\textbf{Answer:} }

\begin{document}
\sheettitle{Sheet 8: From states to observables}{Short-distance quantities, radiative overlaps, and strong decays}{Prerequisite: Sheets 1--7. \difficulty{★★/★★★} \quad \timebox{2--3 hours}}

\begin{goals}
Relate wavefunctions to decay widths; derive phase-space scaling and overlap
selection rules; understand why masses do not validate wavefunctions; analyze
how mixing propagates coherently into amplitudes.
\end{goals}
\begin{reading}
Selected reading in Landau--Lifshitz, \emph{Quantum Mechanics: Non-Relativistic
Theory}: Ch. VI, ``Perturbation Theory,'' for transition probabilities, and
Ch. XIV, ``Addition of Angular Momenta,'' for angular selection rules. This is
optional background; the sheet is not built as a commentary on the book.
\end{reading}


\begin{problem}{Problem 1 ★ — Derive two-body momentum and threshold scaling}
For $A\to B+C$, derive
$k=\sqrt{[M_A^2-(M_B+M_C)^2][M_A^2-(M_B-M_C)^2]}/(2M_A)$.
If the final relative partial wave is $\ell$, show why
$\Gamma\propto k^{2\ell+1}$ near threshold.
\end{problem}
\begin{solution}
In the parent rest frame, energy conservation and on-shell conditions give the
K\"all\'en-function expression above. A partial-wave amplitude behaves as
$k^\ell$ at small $k$, while two-body phase space supplies one more power of
$k$, hence $\Gamma\propto k^{2\ell+1}$.
\end{solution}
\begin{quiz}
Why can a few-MeV mass error cause a large width error near threshold? \answer
the available momentum is small and enters with a positive odd power.
\end{quiz}

\begin{problem}{Problem 2 ★★ — Derive an electric-dipole radial rule}
For the E1 operator $\bm r$, use parity and angular momentum to show
$\Delta L=\pm1$ and no spin flip at leading order. Express the radial factor in
terms of reduced waves.
\end{problem}
\begin{solution}
$\bm r$ is a rank-1 odd-parity tensor, giving the triangle rule
$|L-L'|\le1$ and opposite orbital parity; thus $\Delta L=\pm1$. It does not act
on spin, so $\Delta S=0$. The radial matrix element is
$\int_0^\infty u_f(r)\,r\,u_i(r)\dd r$.
\end{solution}
\begin{quiz}
Can two solvers agree on masses but disagree on an E1 rate? \answer yes;
eigenvalues constrain the Hamiltonian globally, while the rate tests a weighted
wavefunction overlap.
\end{quiz}

\begin{problem}{Problem 3 ★★ — Explain origin sensitivity}
For a normalized $S$ wave show $|\Psi(0)|^2=|R(0)|^2/(4\pi)$ and relate
$R(0)$ to $u'(0)$. Why is using $u(h)/h$ only first-order safe without further
analysis?
\end{problem}
\begin{solution}
$Y_{00}=1/\sqrt{4\pi}$ and regularity gives $u(r)=R(0)r+O(r^2)$, so
$R(0)=u'(0)$. If values are known only at radii separated by $h$, then
$u(h)/h=R(0)+O(h)$ in general. One must study the limit as $h$ decreases, or
evaluate the derivative directly in a basis representation, before treating it
as an accurate value of $R(0)$.
\end{solution}
\begin{quiz}
Which is usually more demanding numerically: a mass or $|\Psi(0)|^2$? \answer
the origin density, because it depends on local derivative information.
\end{quiz}

\begin{problem}{Problem 4 ★★ — Propagate mixing into a decay amplitude}
If $\ket i=\cos\theta\ket a+\sin\theta\ket b$, derive the decay width to a fixed
final channel in terms of amplitudes $\mathcal M_a,\mathcal M_b$. Identify the
interference term.
\end{problem}
\begin{solution}
$\mathcal M_i=\cos\theta\mathcal M_a+\sin\theta\mathcal M_b$, hence
\[
|\mathcal M_i|^2=\cos^2\theta|\mathcal M_a|^2+
\sin^2\theta|\mathcal M_b|^2+2\sin\theta\cos\theta
\operatorname{Re}(\mathcal M_a^*\mathcal M_b).
\]
Adding component widths would omit the last term.
\end{solution}
\begin{quiz}
Why does a basis-phase choice not change the interference physically? \answer
the corresponding component coefficient and matrix element change sign
together.
\end{quiz}

\begin{problem}{Problem 5 ★★★ — Separate spin algebra from wave distortion}
Suppose two annihilation widths have
$\Gamma_0/\Gamma_2=(15/4)\,|D_0/D_2|^2$, where $15/4$ is the short-distance
spin-color factor and $D_J$ is a number obtained by applying a specified linear
integral or derivative to the entire wavefunction. What conclusion follows if
the ratio differs from $15/4$? Design two controlled comparisons.
\end{problem}
\begin{solution}
The deviation need not mean wrong spin algebra; it measures $J$-dependent wave
distortion when $D_0\ne D_2$. First evaluate both widths with a common central
wave, which should isolate $15/4$. Then use separately solved $J$ waves while
holding the short-distance coefficients fixed; the change isolates the radial
effect. This two-step comparison separates spin algebra from radial-wave effects
and is more informative than fitting one overall constant.
\end{solution}
\begin{quiz}
What does a decay calculation validate that a spectrum fit cannot? \answer wave
normalization, local/overlap structure, mixing phases, and operator selection
rules.
\end{quiz}

\begin{paperbridge}
GI Sec. IV tests the calculated states through their couplings: Sec. IV A treats
strong decays, Secs. IV B--C treat electromagnetic annihilation and photon
emission, and Sec. IV D collects further couplings. The principal numerical
results appear in Tables IV--VII. These observables depend on the same masses,
wavefunctions, and mixed-state components constructed in Secs. II--III.
\end{paperbridge}
\end{document}
